% Source: https://artofproblemsolving.com/community/c3h2900424_jmpsc_unused_proposals
% Stable identity: aops-topic-2900424-q10
The circumcircle of such a triangle is the polygon's circumcircle, so the right angle's opposite side must be a diameter. For odd $n$ there are no diameters, and $T_n=0$.
\begin{asy}
usepackage("amsmath");
size(2.5cm);
defaultpen(fontsize(10pt));
pair[] P;for(int i=0;i<8;++i)P[i]=dir(22.5+45*i);draw(P[0]--P[1]--P[2]--P[3]--P[4]--P[5]--P[6]--P[7]--cycle);draw(circle((0,0),1),gray);draw(P[1]--P[5],dashed);draw(P[1]--P[3]--P[5]);
\end{asy}

For $n=2x$ there are $x$ diameters, each completed by any of the other $2x-2$ vertices, so $T_{2x}=2x(x-1)=4\dbinom x2$. By the hockey-stick identity,
\[\sum_{x=2}^{23}4\dbinom x2=4\dbinom{24}3=\boxed{8096}.\]
